\subsection{Момент сил}
%%%%%%%%%%%%%%%%%%%%%%%%%%%%%%%%%%%%%%%%%%%%%%%%%%%%
%%%%%%%%%%%%%%%%%%%%%%%%%%%%%%%%%%%%%%%%%%%%%%%%%%%%
%%%%%%%%%%%%%%%%%%%%%%%%%%%%%%%%%%%%%%%%%%%%%%%%%%%%
\z[{\begin{wrapfigure}[3]{r}{0.25\textwidth}%
		\vspace{-10pt}%
		\begin{flushright}%
 			\includegraphics[width=0.25\textwidth]{5.2.1}%
		\end{flushright}%
	\end{wrapfigure}%
	}]%
{В системе, приведённой на рисунке, масса первого груза равна $m$, масса второго в $a = 2$ раза больше, а масса третьего в $b = 3$ раза меньше. Масса рычага равна $M = 18\,\text{кг}$. Чему равна масса $m$?
}
%
{$m = \frac{3}{38} M = 1.42\,\text{кг}$}

%%%%%%%%%%%%%%%%%%%%%%%%%%%%%%%%%%%%%%%%%%%%%%%%%%%%
\z[{\begin{wrapfigure}[5]{r}{0.25\textwidth}%
		\vspace{+5pt}%
		\begin{flushright}%
 			\includegraphics[width=0.25\textwidth]{5.2.2}%
		\end{flushright}%
	\end{wrapfigure}%
	}]%
{Система, состоящая из двух однородных стержней разной плотности, находится в равновесии. Масса верхнего стержня $m_1 = 1.4\,\text{кг}$. Трение пренебрежимо мало. Определите, при какой массе $m_2$ нижнего стержня возможно такое равновесие.
}
%
{$m_1 = m_2$}

%%%%%%%%%%%%%%%%%%%%%%%%%%%%%%%%%%%%%%%%%%%%%%%%%%%%
\z[{\begin{wrapfigure}[6]{r}{0.25\textwidth}%
		\vspace{-25pt}%
		\begin{flushright}%
 			\includegraphics[width=0.25\textwidth]{5.2.3}%
		\end{flushright}%
	\end{wrapfigure}%
	}]%
{Рычаг, размеченный на 7 одинаковых частей, может свободно вращаться на опоре (см. рисунок). На нём стоит груз массой $m = 2\,\text{кг}$, привязанный с помощью нити, перекинутой через блок к концу рычага. Кроме этого, снизу к рычагу прикреплён груз массой $3m$. При каком максимальном значении массы рычага $M$ возможно равновесие системы?
\vspace{2em}
}
%
{$M = 10m = 20\,\text{кг}$}


%%%%%%%%%%%%%%%%%%%%%%%%%%%%%%%%%%%%%%%%%%%%%%%%%%%%
\z[{\begin{wrapfigure}[3]{r}{0.2\textwidth}%
		\vspace{-30pt}%
		\begin{flushright}%
 			\includegraphics[width=0.2\textwidth]{5.2.4}%
		\end{flushright}%
	\end{wrapfigure}%
	}]%
{Изображённая на рисунке система из рычага и блоков находится в равновесии. Точки подвеса делят рычаг в отношении $a:b$. Найдите отношение масс грузов.
\vspace{2em}
}
%
{$\frac{m_a}{m_b} = \frac{2a + 4b}{3a + b}$}

%%%%%%%%%%%%%%%%%%%%%%%%%%%%%%%%%%%%%%%%%%%%%%%%%%%%
\z[{\begin{wrapfigure}[3]{r}{0.2\textwidth}%
		\vspace{-15pt}%
		\begin{flushright}%
 			\includegraphics[width=0.2\textwidth]{5.2.5}%
		\end{flushright}%
	\end{wrapfigure}%
	}]%
{При каких значениях массы груза $M$ возможно равновесие системы, изображённой на рисунке? Масса $m$ известна.
\vspace{2em}
}
%
{$m < M < 13m$}


%%%%%%%%%%%%%%%%%%%%%%%%%%%%%%%%%%%%%%%%%%%%%%%%%%%%
\z[{\begin{wrapfigure}[8]{r}{0.2\textwidth}%
		\vspace{-15pt}%
		\begin{flushright}%
 			\includegraphics[width=0.2\textwidth]{5.2.6}%
		\end{flushright}%
	\end{wrapfigure}%
	}]%
{Для удержания тяжёлого груза используется система из шести блоков и нескольких тросов, прикреплённых к потолку так, как показано на рисунке. С какой силой $F$ надо тянуть вниз за конец троса, свисающего с левого блока, чтобы удерживать груз массой $M = 500\,\text{кг}$ в равновесии? Участки тросов, не лежащие на блоках, вертикальны; весом блоков и тросов, а также трением можно пренебречь. Ускорение свободного падения принять равным $g = 10\,\text{м/с}^2$.}
%
{$F = \frac{Mg}{10} = 500\,\text{Н}$}

%%%%%%%%%%%%%%%%%%%%%%%%%%%%%%%%%%%%%%%%%%%%%%%%%%%%
\z[{\begin{wrapfigure}[2]{r}{0.35\textwidth}%
		\vspace{-15pt}%
		\begin{flushright}%
 			\includegraphics[width=0.35\textwidth]{5.2.7}%
		\end{flushright}%
	\end{wrapfigure}%
	}]%
{При каких значениях массы рычага $M$ система будет находиться в равновесии.}
%
{$2m<M<3m$}
